\QHeading{第1题　脱式计算。}
\par\noindent


\CalcBlock{1}{\begin{aligned}[t]&\phantom{={}}\color{black}900-25\times24+96\\&\color{answerred}=900-600+96\\&\color{answerred}=300+96\\&\color{answerred}=396\end{aligned}}
\par

\CalcBlock[简便解法（补充）]{2}{\begin{aligned}[t]&\phantom{={}}\color{black}540-36\times 22\div 9\\&\color{answerred}=540-(36\div 9)\times 22\\&\color{answerred}=540-4\times 22\\&\color{answerred}=540-88\\&\color{answerred}=452\end{aligned}}
\MethodBlock{按运算顺序计算（基础练习）}{本单元还要练习运算顺序；与上面的简便解法对照。}{\begin{aligned}[t]&\phantom{={}}540-36\times22\div9\\&=540-792\div9\\&=540-88\\&=452\end{aligned}}



\CalcBlock{3}{\begin{aligned}[t]&\phantom{={}}\color{black}25\times50-30\times20\\&\color{answerred}=1250-600\\&\color{answerred}=650\end{aligned}}
\CalcBlock[简便解法（补充）]{4}{\begin{aligned}[t]&\phantom{={}}\color{black}860-42\times 13+40\\&\color{answerred}=860-546+40\\&\color{answerred}=(860+40)-546\\&\color{answerred}=900-546\\&\color{answerred}=354\end{aligned}}
\MethodBlock{按运算顺序计算（基础练习）}{本单元还要练习运算顺序；与上面的简便解法对照。}{\begin{aligned}[t]&\phantom{={}}860-42\times13+40\\&=860-546+40\\&=314+40\\&=354\end{aligned}}



\CalcBlock[方法一]{5}{\begin{aligned}[t]&\phantom{={}}\color{black}100\div5+120\div6\times3\\&\color{answerred}=20+20\times3\\&\color{answerred}=20+60\\&\color{answerred}=80\end{aligned}}
\MethodBlock{方法二：先处理括号或乘除，再提取相同因数}{用乘法分配律反过来计算。}{\begin{aligned}[t]&\phantom{={}}100\div 5+120\div 6\times 3\\
&=20+20\times 3\\
&=20\times (1+3)\\
&=20\times 4\\
&=80\end{aligned}}

\CalcBlock{6}{\begin{aligned}[t]&\phantom{={}}\color{black}200-60\div5\times9\\&\color{answerred}=200-12\times9\\&\color{answerred}=200-108\\&\color{answerred}=92\end{aligned}}
\par\noindent

\QHeading{第2题　火眼金睛辨对错。}

\par\noindent\makebox[\linewidth][c]{
\begin{tikzpicture}[x=1bp,y=-1bp,color=black,draw=black,text=black]
\begin{scope}\clip (56.0,415.0) rectangle (549.0,541.33);
% 高清原题重排：PDF 7，原书 2。不含解答。
\hdhead{2}{四则混合运算——信息窗1}
\hdsection{1}{148}{脱式计算。}
\q{100}{193}{245}{13.5}{\(900-25\times24+96\)}
\q{355}{193}{245}{13.5}{\(540-36\times22\div9\)}
\q{611}{193}{245}{13.5}{\(25\times50-30\times20\)}
\q{100}{383}{245}{13.5}{\(860-42\times13+40\)}
\q{355}{383}{245}{13.5}{\(100\div5+120\div6\times3\)}
\q{611}{383}{245}{13.5}{\(200-60\div5\times9\)}
\hdsection{2}{588}{火眼金睛辨对错。}
\q{103}{632}{34}{13.5}{①}
\q{461}{632}{34}{13.5}{②}
\q{157}{631}{171}{13.5}{\(240-40\div5\times9\)}
\q{335}{631}{105}{13.5}{改正：}
\q{139}{675}{190}{13.5}{\(=200\div5\times9\)}
\q{139}{719}{185}{13.5}{\(=40\times9\)}
\q{139}{761}{180}{13.5}{\(=360(\qquad)\)}
\q{516}{631}{190}{13.5}{\(120\times2+150\div3\)}
\q{707}{631}{105}{13.5}{改正：}
\q{495}{675}{190}{13.5}{\(=240+150\div3\)}
\q{495}{719}{190}{13.5}{\(=390\div3\)}
\q{495}{761}{180}{13.5}{\(=130(\qquad)\)}
\hdsection{3}{815}{解决问题。}
\q{104}{861}{745}{13.5}{王叔叔建了两个温室大棚分别种植草莓和樱桃，其中种植草莓的温室}
\q{104}{903}{745}{13.5}{大棚占地310平方米，平均每平方米可采摘4千克草莓；种植樱桃的温}
\q{104}{945}{745}{13.5}{室大棚占地198平方米，平均每平方米可采摘11千克樱桃。这两个温}
\q{104}{987}{745}{13.5}{室大棚一共可采摘草莓和樱桃多少千克？}

\begin{scope}[color=answerred,draw=answerred,text=answerred]\ans{66.67}{154.00}{163.33}{11}{\(\begin{aligned}[t]&=900-600+96\\&=300+96\\&=396\end{aligned}\)}
\ans{237.33}{154.00}{161.33}{11}{\(\begin{aligned}[t]&=540-792\div9\\&=540-88\\&=452\end{aligned}\)}
\ans{408.00}{154.00}{162.00}{11}{\(\begin{aligned}[t]&=1250-600\\&=650\end{aligned}\)}
\ans{66.67}{280.67}{163.33}{11}{\(\begin{aligned}[t]&=860-546+40\\&=314+40\\&=354\end{aligned}\)}
\ans{237.33}{280.67}{165.33}{11}{\(\begin{aligned}[t]&=20+20\times3\\&=20+60\\&=80\end{aligned}\)}
\ans{408.00}{280.67}{161.33}{11}{\(\begin{aligned}[t]&=200-12\times9\\&=200-108\\&=92\end{aligned}\)}
\node[anchor=center,inner sep=0pt,outer sep=0pt,text=answerred,font=\normalsize] at (151.333,512.667) {\(\times\)};
\ans{218.00}{448.00}{102.00}{9}{\linespread{1}\selectfont\setlength{\jot}{-2pt}\(\begin{aligned}[t]&\phantom{={}}240-40\div5\times9\\&=240-8\times9\\&=240-72\\&=168\end{aligned}\)}
\node[anchor=center,inner sep=0pt,outer sep=0pt,text=answerred,font=\normalsize] at (390.667,512.667) {\(\times\)};
\ans{464.00}{447.33}{106.00}{9}{\linespread{1}\selectfont\setlength{\jot}{-2pt}\(\begin{aligned}[t]&\phantom{={}}120\times2+150\div3\\&=240+50\\&=290\end{aligned}\)}
\ans{70.00}{691.33}{487.33}{11}{\linespread{1}\selectfont\setlength{\jot}{-2pt}\(\begin{aligned}[t]&\phantom{={}}4\times310+11\times198\\&=1240+2178\\&=3418\text{（千克）}\end{aligned}\)\\答：一共可采摘草莓和樱桃3418千克。}
\end{scope}
\end{scope}\pgfresetboundingbox\path[use as bounding box] (56.0,415.0) rectangle (549.0,541.33);
\end{tikzpicture}}\par



\par



\par\noindent

\QHeading{第3题　解决问题。}

\Problem{王叔叔建了两个温室大棚分别种植草莓和樱桃，其中种植草莓的温室大棚占地310平方米，平均每平方米可采摘4千克草莓；种植樱桃的温室大棚占地198平方米，平均每平方米可采摘11千克樱桃。这两个温室大棚一共可采摘草莓和樱桃多少千克？}
\par


\MethodTitle{方法一}
\Answer{\Formula{\begin{aligned}[t]&\phantom{={}}4\times310+11\times198\\&=1240+2178\\&=3418\text{（千克）}\end{aligned}}\par 答：一共可采摘草莓和樱桃3418千克。}

\MethodTitle{方法二：统一每平方米的采摘量，再补足}
\Answer{先把两棚都按每平方米4千克计算，再补上樱桃棚每平方米少算的7千克。}
\Answer{\Formula{\begin{aligned}[t]&\phantom{={}}4\times (310+198)+(11-4)\times 198\\
&=4\times 508+7\times 198\\
&=2032+1386\\
&=3418\text{（千克）}\end{aligned}}}
\Answer{答：一共可采摘草莓和樱桃3418千克。}
